\documentclass{article}

\usepackage{amsmath, amssymb, bm}
\usepackage[a4paper, margin=1in]{geometry}

\title{A simple AHRS navigation system}
\author{}
\date{}

\begin{document}
	
	\maketitle
	
	
\section{Introduction}

	An \textit{Attitude and Heading Reference System} (AHRS) is the system that will estimate the orientation of our drone in three-dimensional space. An AHRS outputs of roll, pitch, and heading (yaw), which describe the attitude of the body relative to a reference frame.
	
	\vspace{1em}
	
	\noindent An AHRS typically combines data from gyroscopes, accelerometers, and magnetometers. Through sensor fusion algorithms, these measurements are used to produce a stable and accurate estimate of orientation while mitigating the limitations of individual sensors.
	
\section{Reference Frame Definitions}

Two reference frames are used throughout this work: the body-fixed frame ($b$) and the world-fixed frame ($n$).
The body-fixed frame ($b$) is attached to the drone and moves with it. Its axes are defined as follows:

\begin{itemize}
	\item The x-axis points along the forward direction of the drone. Rotation about this axis is referred to as \emph{roll}.
	\item The y-axis points to the right side of the drone. Rotation about this axis is referred to as \emph{pitch}.
	\item The z-axis points upward. Rotation about this axis is referred to as \emph{yaw}.
	\item The coordinate system is right-handed.
\end{itemize}

\noindent The world-fixed frame ($n$) is an inertial reference frame following the East--North--Up (ENU) convention:

\begin{itemize}
	\item The x-axis points toward the East.
	\item The y-axis points toward the North.
	\item The z-axis points upward.
	\item The coordinate system is right-handed.
\end{itemize}


	
\section{The conceptual idea}

	In this section, we present the conceptual foundation of our AHRS implementation. The system relies on three sensors:

	\begin{itemize}
		\item Gyroscope: measures the angular velocity vector $\boldsymbol{\omega}$ around each axis of the body-fixed frame $b$.
		
		\item Accelerometer: measures the specific force acting on the body as a three-dimensional vector in the body-fixed frame $b$. When the body is not subject to linear acceleration, the measurement corresponds to the gravity vector.
		
		\item Magnetometer: measures the local Earth's magnetic field as a three-dimensional vector in the body-fixed frame $b$. This measurement provides a reference for determining heading relative to magnetic north.

		
	\end{itemize}

	\noindent The gyroscope provides the high-frequency component of the orientation estimate since the measured angular velocity can be integrated to determine the body's orientation over time. However, relying solely on gyroscope measurements is not sufficient, since sensor biases and noise accumulate during integration, causing the estimated orientation to drift over time.
	
	\vspace{1em}
	
	\noindent To mitigate this drift, measurements from the accelerometer and magnetometer are incorporated. The accelerometer cannot directly provide the gravity vector during flight, since its measurements are corrupted by the drone's linear accelerations. Nevertheless, under the assumption that these accelerations average to zero over sufficiently long time intervals, the average accelerometer measurement converges to the gravity vector. Similarly, the magnetometer provides a measurement of the Earth's magnetic field, which is assumed to be constant in the earth-fixed frame $n$.
	
	\vspace{1em}
	
	\noindent Consequently, the gravity and magnetic field vectors serve as stable low-frequency references that can be used to correct the long-term drift of the orientation estimate while preserving the gyroscope's fast dynamic response.
	
\section{Description of the Algorithm}

\subsection{Computing New Orientation}

Mathematically, we describe the orientation of the drone as a matrix $R_{k}^{nb}$. This matrix transforms 3D points or vectors described in the world-fixed $n$ coordinate system to the body-fixed $b$ coordinate system. The inverse transformation ($b \rightarrow n$) can be done by simply inverting (or transposing, since it is orthogonal) the matrix:

\[
R^{bn}_k = (R^{nb}_k)^{-1} = (R^{nb}_k)^{T}
\]

\noindent We can use the elements in this $3\times3$ matrix to derive the Euler angles (roll, pitch, yaw).
\vspace{1em}

\noindent Given the orientation estimate from the previous time step, denoted by $R_k^{nb}$, the orientation at the next time step can be predicted using the angular velocity measurements provided by the gyroscope. The gyroscope measures the instantaneous angular velocity of the drone expressed in the body-fixed frame $b$.

\vspace{1em}
\noindent Assuming that the angular velocity remains constant during the sampling interval ($\Delta t$), the orientation can be propagated forward in time. Because the angular velocity $\boldsymbol{\omega}$ is measured relative to the local body frame, the rotation update must pre-multiply our orientation matrix:

\[
R_{k+1}^{nb} = e^{\Omega(\boldsymbol{\omega})\Delta t} R_{k}^{nb}
\]

\noindent To make this equation suitable for real-time computation on an embedded system, we can expand the matrix exponential using a Taylor series approximation:

\[
e^{\Omega(\boldsymbol{\omega})\Delta t} = I + \Omega(\boldsymbol{\omega})\Delta t + \frac{1}{2!}\left(\Omega(\boldsymbol{\omega})\Delta t\right)^2 + \frac{1}{3!}\left(\Omega(\boldsymbol{\omega})\Delta t\right)^3 + \dots
\]

\noindent Given that the sampling interval is highly frequent (e.g., $\Delta t = 0.01\,\mathrm{s}$), the higher-order terms ($(\Delta t)^2, (\Delta t)^3, \dots$) become negligibly small. Truncating the series to a first-order linear approximation yields:

\[
e^{\Omega(\boldsymbol{\omega})\Delta t} \approx I + \Omega(\boldsymbol{\omega})\Delta t
\]

\noindent Substituting this back into the propagation step results in a computationally lightweight update equation:

\[
R_{k+1}^{nb} \approx \left(I + \Omega(\boldsymbol{\omega})\Delta t\right) R_{k}^{nb}
\]

\noindent Where:
\begin{itemize}
	\item $R_{k+1}^{nb}$ is a $3\times3$ rotation matrix describing the new uncorrected orientation at time step $k+1$.
	\item $I$ is the $3\times3$ identity matrix.
	\item $\Omega(\boldsymbol{\omega})$ is the skew-symmetric $3\times3$ matrix encoding the body-frame angular velocity vector $\boldsymbol{\omega} = [\omega_x, \omega_y, \omega_z]^T$:
	\[
	\Omega(\boldsymbol{\omega}) = \begin{bmatrix} 
		0 & -\omega_z & \omega_y \\ 
		\omega_z & 0 & -\omega_x \\ 
		-\omega_y & \omega_x & 0 
	\end{bmatrix}
	\]
\end{itemize}

\subsection{Estimation of Reference Vectors}

As discussed previously, the accelerometer and magnetometer are used to estimate low-frequency reference vectors that can be used to correct the drift that accumulates in the gyroscope-based orientation estimate over time.

\vspace{1em}
\noindent For the accelerometer, we know that when the drone is not accelerating, the measured acceleration corresponds solely to the gravitational acceleration vector. During flight, however, the accelerometer also measures accelerations caused by vehicle motion. We assume that these additional accelerations have a zero mean over a sufficiently long time interval. Consequently, applying a low-pass filter to the accelerometer measurements yields an estimate of the gravity vector. The same principle can be applied to the magnetometer measurements to obtain an estimate of the Earth's magnetic field vector.

\vspace{1em}
\noindent The following procedure describes the estimation of the gravity vector. The estimation of the magnetic north vector is performed analogously.

\vspace{1em}
\noindent First, the measured acceleration vector ($a_m^b$), expressed in the body-fixed frame ($b$), is transformed into the navigation frame ($n$):

\[
a_m^n = \left(R_k^{nb}\right)^{-1} a_m^b.
\]

\vspace{1em}
\noindent Next, an exponential moving average is updated using the transformed measurement and the previous estimate ($a_k^n$):

\[
a_{k+1}^n = \alpha a_m^n + (1-\alpha)a_k^n.
\]

\vspace{1em}
\noindent The filter coefficient ($\alpha$) is chosen as

\[
\alpha = \frac{2}{N+1}
\]

\noindent where (N) is the effective averaging window expressed in samples.
For example, if an averaging window of 60 seconds is desired and the sampling period is ($\Delta t = 0.01,\mathrm{s}$), then

\[
N = \frac{60}{0.01} = 6000
\]

\noindent As ($N$) increases, the estimate becomes less sensitive to short-term disturbances and converges more closely to the underlying gravity vector.

\vspace{1em}
\noindent Finally we need to transform the new estimate back to the $b$ frame:
\[
g^b_m = R_k^{nb} a_{k+1}^n
\]
We can now use the $g^b_m$ as an estimate of the gravity vector.
A similar process can be used to obtain an a vector $m^b_m$ pointing to the magnetic north.


\subsection{Drift Correction}

\noindent In this section, we discuss how to correct for the drift caused by integrating the angular velocities of the gyroscope.

\vspace{1em}
\noindent We use a measured estimate of the gravity vector $g_m^b$ to correct for roll and pitch drift. Given our current estimate of the orientation $R^{nb}_k$, we can compute the theoretical orientation of the gravity vector in the body frame, denoted as $g_t^b$, and then compute the angle between them.

\vspace{1em}
\noindent The theoretical orientation of gravity in the navigation ($n$) frame is defined using the acceleration due to gravity:
\[
g^n = 
\begin{bmatrix}
	0 \\
	0 \\
	-9.81
\end{bmatrix}
\]

\noindent We transform this vector into the body ($b$) frame to get our theoretical gravity vector $g_t^b$:
\[
g_t^b = R^{nb}_{k} \cdot g^n
\]

\noindent Next, we compute the error vector between the measured gravity vector and the theoretical one. Because these vectors contain magnitude information, we normalize both to unit vectors before computing their cross product:
\[
r^b = \frac{g_t^b}{\|g_t^b\|} \times \frac{g_m^b}{\|g_m^b\|} 
\]

\noindent From the properties of the cross product of two unit vectors, the magnitude of $r^b$ provides the sine of the angle $\theta$ between the measured gravity vector and the theoretical one, while its direction provides the rotation axis. Therefore, assuming the error angle is small, we can state that:
\[
\|r^b\| = \sin \theta \approx \theta
\]

\noindent We can extract the normalized axis of rotation $u^b$ by isolating the direction of the error vector:
\[
u^b = \frac{r^b}{\|r^b\|}
\]

\noindent To convert an arbitrary axis-angle representation into a rotation matrix, we can utilize Rodrigues' rotation formula:
\[
R = I + \sin\theta\,[u^b]_\times + (1-\cos\theta)\,[u^b]_\times^2
\]

\noindent Please note $[u^b]_\times$ is the skew symmetric form of the unit vector $u$.

\vspace{1em}
\noindent In our specific case, we do not wish to correct for the entire angle $\theta$ at once. We assume that the high-frequency gyroscope data is more accurate in the short term, and we use this gravity correction primarily to apply a small, continuous offset that compensates for long-term drift. Therefore, we only compensate for a fraction of $\theta$. We define this correction gain as $K_g$, where a typical value is $K_g = 0.01$.

\vspace{1em}
\noindent The actual correction angle we wish to apply is:
\[
\alpha = K_g \cdot \theta
\]

\noindent Assuming $\alpha$ is a small angle, we can use the small-angle approximations $\sin\alpha \approx \alpha$ and $\cos\alpha \approx 1$. Applying these approximations to Rodrigues' formula for the correction angle $\alpha$ and the rotation axis $u^b$ allows us to simplify it to a linear approximation.
\[
R_g \approx I + \alpha \cdot [u^b]_\times
\]

\noindent We can use an identical process to obtain the correction matrix for the yaw direction $R_m$ using the magnetometer. The only difference being that we start with the magnetic north direction $m^n$ and use the reference vector $m^b_m$:
\[
m^n = 
\begin{bmatrix}
	0 \\
	1 \\
	0
\end{bmatrix}
\]

\noindent We can now obtain the final orientation by applying the two drift correction terms:
\[
	R = R_g \cdot R_m \cdot R_{k+1} 
\]

\subsection{Normalization of the Rotation Matrix}

The first-order integration method used to propagate the orientation,

\[
R_{k+1}^{nb} \approx \left(I + \Omega(\boldsymbol{\omega})\Delta t\right)R_k^{nb},
\]

does not exactly preserve the orthogonality of the rotation matrix. Due to numerical approximation errors, the rows and columns of the matrix gradually lose their unit length and mutual orthogonality. If left uncorrected, the matrix will eventually cease to represent a valid rotation.

\vspace{1em}
\noindent To maintain a valid orientation estimate, the rotation matrix must be periodically re-orthogonalized. A computationally efficient method is based on the Gram--Schmidt process.

\vspace{1em}
\noindent Let the rows of the estimated rotation matrix be

\[
R^{nb} =
\begin{bmatrix}
	r_1^T \\
	r_2^T \\
	r_3^T
\end{bmatrix}.
\]

First, the first row is normalized:

\[
\hat r_1 = \frac{r_1}{\|r_1\|}.
\]

Next, the component of the second row along the first row is removed:

\[
r_2' = r_2 - (\hat r_1 \cdot r_2)\hat r_1,
\]

and the result is normalized:

\[
\hat r_2 = \frac{r_2'}{\|r_2'\|}.
\]

Finally, the third row is reconstructed using the cross product:

\[
\hat r_3 = \hat r_1 \times \hat r_2.
\]

The corrected rotation matrix is then

\[
R^{nb} =
\begin{bmatrix}
	\hat r_1^T \\
	\hat r_2^T \\
	\hat r_3^T
\end{bmatrix}.
\]

\vspace{1em}
\noindent This procedure guarantees that the rows remain mutually orthogonal and of unit length, thereby ensuring that the matrix remains a valid rotation matrix satisfying

\[
(R^{nb})^T R^{nb} = I.
\]

\vspace{1em}
\noindent In practice, this normalization step should be performed after each orientation update cycle to prevent the accumulation of numerical errors.

\subsection{Extraction of Euler Angles}

Once the orientation matrix $R_{k}^{nb}$ is updated and corrected, the vehicle's attitude can be expressed in terms of traditional Euler angles: roll ($\phi$), pitch ($\theta$), and yaw ($\psi$). Following the standard aerospace convention for an East--North--Up (ENU) navigation frame transforming into a body frame ($n \rightarrow b$), the matrix is constructed via a sequential Yaw-Pitch-Roll ($Z$-$Y$-$X$) rotation sequence:

\[
R^{nb}(\phi, \theta, \psi) = R_x(\phi) R_y(\theta) R_z(\psi)
\]

\noindent Evaluating this product yields the algebraic structure of our $3\times3$ orientation matrix in terms of the individual angles:

\[
R^{nb} = \begin{bmatrix}
	r_{11} & r_{12} & r_{13} \\
	r_{21} & r_{22} & r_{23} \\
	r_{31} & r_{32} & r_{33}
\end{bmatrix} = 
\begin{bmatrix}
	\cos\theta\cos\psi & \cos\theta\sin\psi & -\sin\theta \\
	\sin\phi\sin\theta\cos\psi - \cos\phi\sin\psi & \sin\phi\sin\theta\sin\psi + \cos\phi\cos\psi & \sin\phi\cos\theta \\
	\cos\phi\sin\theta\cos\psi + \sin\phi\sin\psi & \cos\phi\sin\theta\sin\psi - \sin\phi\cos\psi & \cos\phi\cos\theta
\end{bmatrix}
\]

\noindent Under normal flight dynamics where the drone does not experience extreme vertical attitudes, the Euler angles can be isolated directly using the multi-quadrant inverse tangent function ($\text{atan2}$) and the inverse sine function:

\begin{align*}
	\phi   &= \text{atan2}\left(r_{23}, r_{33}\right) \\
	\theta &= -\arcsin\left(r_{13}\right) \\
	\psi   &= \text{atan2}\left(r_{12}, r_{11}\right)
\end{align*}

\subsubsection{Handling Gimbal Lock}

A mathematical singularity known as \emph{Gimbal Lock} occurs if the drone pitches straight up ($\theta = 90^\circ$) or straight down ($\theta = -90^\circ$). In these states, $\cos\theta = 0$, causing the terms $r_{11}, r_{12}, r_{23},$ and $r_{33}$ to become zero. This collapses roll and yaw onto the same geometric axis, making it impossible to solve for them independently. To prevent numerical instability in the firmware, conditional branching must be applied:

\vspace{1em}
\noindent \textbf{Case 1: Pitching Straight Up ($r_{13} \approx -1$)} \\
The pitch angle is locked to $\theta = \pi/2$. We arbitrarily set the roll angle to zero ($\phi = 0$) to resolve the degree of freedom, allowing yaw to be cleanly computed as:
\[
\psi = \text{atan2}\left(r_{21}, r_{22}\right)
\]

\vspace{1em}
\noindent \textbf{Case 2: Pitching Straight Down ($r_{13} \approx 1$)} \\
The pitch angle is locked to $\theta = -\pi/2$. Setting the roll angle to zero ($\phi = 0$), yaw is computed as:
\[
\psi = -\text{atan2}\left(r_{21}, r_{22}\right)
\]

\subsection{System Initialization}

Before the AHRS can reliably track orientation using gyroscope integration, an initial orientation matrix $R_0^{bg}$ must be established. When the drone is powered on and stationary on the ground, the initial acceleration measurement corresponds purely to gravity, and the magnetometer measures the static local Earth's magnetic field. We can exploit these physical vector readings to construct a deterministic, orthogonal coordinate system via the Triad method.

A rotation matrix $R^{sd}$ transforms a vector from a source frame $s$ to a destination frame $d$. Conceptually, it is structured such that \textbf{each column represents a unit basis vector of the source frame $s$ expressed in the coordinates of the destination frame $d$.} Because our sensor measurements naturally provide the navigation frame's axes (gravity and magnetic north) expressed in body-frame coordinates, it is most direct to first construct the matrix $R_0^{gb}$ (from navigation $g$ to body $b$). The columns of $R_0^{gb}$ will simply be the navigation frame's basis vectors (East, North, Up) expressed in body-frame coordinates.

\vspace{1em}
\noindent Let $a_{\text{init}}^b$ and $m_{\text{init}}^b$ be the initial normalized vector measurements of the accelerometer and magnetometer expressed in the body frame:
\[
a^b = \frac{a_{\text{init}}^b}{\|a_{\text{init}}^b\|}, \quad m^b = \frac{m_{\text{init}}^b}{\|m_{\text{init}}^b\|}
\]

\noindent According to the ENU reference frame definitions, the true gravity vector points along the negative $z$-axis ($g$-frame), meaning the Up axis ($z$) points directly opposite to gravity. Crucially, a stationary accelerometer measures \textit{proper acceleration}—the upward reaction force exerted by the ground to oppose gravity. Consequently, the positive body-frame acceleration vector $a^b$ points physically upward, meaning it directly represents the tracking of the navigation frame's Up basis vector ($\mathbf{u}^b$) without requiring negation:
\[
\mathbf{u}^b = a^b
\]

\noindent Next, the raw magnetometer vector $m^b$ cannot be used directly as the North basis vector because it is generally not perfectly orthogonal to the gravity vector due to the magnetic inclination angle (dip angle). To extract a clean East-West lateral axis that is strictly perpendicular to the vertical axis, we compute the cross product of the magnetic vector and the vertical axis. This isolates the navigation frame's East basis vector ($\mathbf{e}^b$) expressed in the body frame:
\[
\mathbf{e}^b = \frac{m^b \times \mathbf{u}^b}{\|m^b \times \mathbf{u}^b\|}
\]

\noindent Finally, to complete the right-handed coordinate system and obtain the true horizontal North basis vector ($\mathbf{n}^b$) expressed in the body frame, we take the cross product of our newly computed Up and East basis vectors:
\[
\mathbf{n}^b = \mathbf{u}^b \times \mathbf{e}^b
\]

\noindent Because $\mathbf{e}^b$, $\mathbf{n}^b$, and $\mathbf{u}^b$ are the mutually orthogonal unit basis vectors of the source frame ($g$) expressed in the destination frame ($b$), we pack them directly into the columns of $R_0^{gb}$:
\[
R_0^{gb} = 
\begin{bmatrix}
	\mathbf{e}^b & \mathbf{n}^b & \mathbf{u}^b
\end{bmatrix}
\]

\noindent To obtain the final required initial orientation matrix $R_0^{bg}$, which defines the transformation from the body frame to the navigation frame, we exploit the orthogonal property of rotation matrices ($R^{-1} = R^T$). Transposing $R_0^{gb}$ yields our final row-stacked matrix:
\[
R_0^{bg} = (R_0^{gb})^T = 
\begin{bmatrix}
	(\mathbf{e}^b)^T \\
	(\mathbf{n}^b)^T \\
	(\mathbf{u}^b)^T
\end{bmatrix}
\]

\noindent This matrix maps the initial physical state of the vehicle perfectly onto the navigation frame, establishing a stable baseline for the subsequent integration updates ($R_k^{bg}$).


\end{document}